% Part: applied-modal-logics
% Chapter: epistemic-logic
% Section: bisimulations

\documentclass[../../../include/open-logic-section]{subfiles}

\begin{document}

\olfileid{aml}{el}{bsd}

\olsection{Bisimulasi}

Isih ana pitakonan bab daya ekspresif basa epistemik kita, yaiku
sesambungan antarane model lan rumus kang bener ing model iku.
Saka asil korespondensi kerangka, kita wis weruh yen rumus tartamtu
kang valid ing sawijining kerangka mesthekake yen kerangka iku nduweni
sipat tartamtu. Nanging, apa basa modal kita, contone, bisa mbedakake
antarane jagad kang nduweni panah refleksif lan rerangken jagad tanpa
wates kang saben jagade nyambung menyang jagad sabanjure? Tegese,
apa ana rumus~$A$ kang mung bener ing salah siji saka rong jagad iku?

Bisimulasi iku sesambungan kang bisa kita tetepake antarane model
relasional kanggo nyatakake yen struktur model-model iku padha
miturut pangertèn modal. Kaya kang bakal kita deleng, bisimulasi
ngandharake ekuivalensi antarane model miturut apa kang bisa
dibedakake dening basa epistemik kita.

\begin{defn}[Bisimulasi]
Upamane $M_1 = \tuple{W_1, R_1, V_1}$ lan $M_2 = \tuple{W_2, R_2, V_2}$
iku rong model relasional. Upamane uga $\mathcal{R} \subseteq W_1 \times W_2$
iku relasi biner. Kita nyebut $\mathcal{R}$ minangka
\emph{bisimulasi} yen kanggo saben $\tuple{w_1, w_2} \in \mathcal{R}$,
syarat-syarat iki kawujud:

\begin{enumerate}
  \item $w_1 \in V_1(p)$ yen lan mung yen $w_2 \in V_2(p)$ kanggo saben
  !!{propositional variable}s~$p$.
  \item Kanggo saben agen $a \in G$ lan jagad $v_1 \in W_1$, yen
  $R_{1,a} w_1 v_1$, ana $v_2 \in W_2$ kang nyukupi
  $R_{2,a} w_2 v_2$ lan $\tuple{v_1, v_2} \in
  \mathcal{R}$.
  \item Kanggo saben agen $a \in G$ lan jagad $v_2 \in W_2$, yen
   $R_{2,a} w_2 v_2$, ana $v_1 \in W_1$ kang nyukupi
   $R_{1,a} w_1 v_1$ lan $\tuple{v_1, v_2} \in
   \mathcal{R}$.
\end{enumerate}

Yen ana bisimulasi antarane $M_1$ lan $M_2$ kang nggandhengake jagad
$w_1$ lan $w_2$, kita uga bisa nulis $\tuple{M_1, w_1} \leftrightarroweq
\tuple{M_2, w_2}$, lan nyebut $\tuple{M_1, w_1}$ lan $\tuple{M_2, w_2}$
\emph{bisimilar}.
\end{defn}

Saben klausa ing relasi bisimulasi njamin prakara kang beda-beda.
Klausa 1 njamin yen jagad-jagad bisimilar nyukupi rumus tanpa
modalitas kang padha, amarga klausa iku njamin kecocokan tumrap kabeh
!!{propositional variable}s. Rong klausa liyane, kang sok diarani
``maju'' lan ``mundur,'' njamin yen relasi aksesibilitase nduweni
struktur kang padha.

\begin{thm}
Yen $\tuple{M_1, w_1} \leftrightarroweq \tuple{M_2, w_2}$, mula kanggo
saben !!{formula}~$!A$, $\mSat{M_1}{!A}[w_1]$ yen lan mung yen
$\mSat{M_2}{!A}[w_2]$.
\end{thm}

\begin{figure}
  \begin{center}
    \begin{tikzpicture}[modal]
      \node[world] (w1) {$w_1$}; 
      \node[world] (w2) [above left=of w1]{$w_2$}; 
      \node[world] (w3) [above right=of w1] {$w_3$};
      \draw[<->] (w1) to node {$a$} (w2);
      \draw[->,loop below] (w1) to node {$a$} (w1);
      \draw[<->] (w1) to [swap] node {$a$} (w3);
      \draw[->,loop above] (w2) to node {$a$} (w2) ;
      \draw[->,loop above] (w3) to node {$a$} (w3) ;
      
      \node[world](v1)[right of=w1, xshift=1.5in]{$v_1$};
      \node[world](v2)[above of=v1]{$v_2$};
      \draw[<->] (v1) to node {$a$} (v2);
      \draw[->,loop below] (v1) to node {$a$} (v1);
      \draw[->,loop above] (v2) to node {$a$} (v2) ;

      \draw[-,dotted] (w1) to (v1) ;
      \draw[-,dotted,bend left=45] (w2) to (v2) ;
      \draw[-,dotted,bend left=30] (w3) to (v2) ;
    \end{tikzpicture}
  \end{center}
  \caption{Rong model kang bisimilar yen valuasine cocog.}
  \ollabel{fig:bisimilar}
\end{figure}

Sanadyan rong model ing \olref{fig:bisimilar} ora persis padha,
yen valuasi jagad-jagad kang digandhengake cocog, ana bisimulasi
kang nggandhengake $w_1$ karo~$v_1$. Bisimulasi iki uga nggandhengake
$w_2$ lan $w_3$ karo~$v_2$: basa modal kita ora bisa medharake
apa-apa kang temenan mbedakake jagad-jagad iku. Nanging, kahanane
bakal beda yen $w_2$ lan~$w_3$ nyukupi !!{propositional variable}s
kang beda.

\end{document}
